% W4i49_Theory.tex
% DFD 17-Agent Research Programme --- Iteration 49 (i49)
% Agents 1 (Alpha), 2 (Topology/k_max), 3 (Fermion Masses), 4 (Spin^c/Exponents), 5 (Anomaly/k-selection)
% Theme: Final theoretical synthesis, programme handoff preparation, Phase 0A requirements
% Date: 2026-03-23

\documentclass[11pt,a4paper]{article}
\usepackage[utf8]{inputenc}
\usepackage{amsmath,amssymb,amsfonts,amsthm}
\usepackage{hyperref}
\usepackage{booktabs}
\usepackage{longtable}
\usepackage{geometry}
\usepackage{xcolor}
\usepackage{tcolorbox}
\geometry{margin=2.5cm}

\newtheorem{theorem}{Theorem}
\newtheorem{proposition}{Proposition}
\newtheorem{lemma}{Lemma}
\newtheorem{corollary}{Corollary}
\newtheorem{definition}{Definition}
\newtheorem{remark}{Remark}

\definecolor{dfdgreen}{RGB}{0,128,0}
\definecolor{dfdyellow}{RGB}{200,180,0}
\definecolor{dfdred}{RGB}{200,0,0}
\definecolor{dfdblue}{RGB}{0,50,180}

\newcommand{\GREEN}{\textcolor{dfdgreen}{\textbf{GREEN}}}
\newcommand{\YELLOW}{\textcolor{dfdyellow}{\textbf{YELLOW}}}
\newcommand{\RED}{\textcolor{dfdred}{\textbf{RED}}}
\newcommand{\NEW}{\textcolor{dfdblue}{\textbf{NEW}}}
\newcommand{\RESOLVED}{\textcolor{dfdgreen}{\textbf{RESOLVED}}}
\newcommand{\Tr}{\operatorname{Tr}}
\newcommand{\CP}{\mathbb{C}P}

\title{DFD i49: Theory --- Alpha, Topology, Masses, Spin$^c$, Anomaly\\[6pt]
\large Agents 1, 2, 3, 4, 5\\[6pt]
\normalsize Iteration 18/20 of Quantum Gravity Cycle (i32--i51)\\[6pt]
\textbf{PENULTIMATE ITERATION --- Final Theoretical Synthesis}}
\author{DFD 17-Agent Research Programme}
\date{2026-03-23}

\begin{document}
\maketitle
\tableofcontents
\newpage

%=============================================================================
\part{Penultimate Iteration: Theory Sector Status Overview}
\label{part:overview}
%=============================================================================

\section{Where Theory Stands at i49}

\begin{tcolorbox}[colback=blue!5!white, colframe=blue!60!black,
  title={Theory Sector: Final Status Entering Penultimate Iteration}]

The DFD theory programme has reached a mature, stable state. The core achievements:

\begin{enumerate}
\item \textbf{Fine structure constant}: $\alpha^{-1} = 137.035\,999\,854$ from $\CP^2 \times S^3$ Chern--Simons level counting with $k_{\max} = 60$. Zero free parameters, sub-ppm precision (+0.005 ppm from CODATA). \GREEN.

\item \textbf{Axiom structure}: 3 independent axioms + 1 derived theorem ($\mu(x) = x/(1+x)$ from $S^3$ composition law). Single free parameter $\lambda$. Prediction-to-parameter ratio 21:1 (honest count).

\item \textbf{VEV formula}: $v = M_P \cdot \alpha^8 \cdot \sqrt{2\pi} = 246.09$ GeV. Confirmed i43 (i42 flag was arithmetic error). Uses standard Planck mass $M_P = 1.221 \times 10^{19}$ GeV. Grade A.

\item \textbf{Alpha\textsuperscript{57} derivation}: Step 9 (modulus stabilization) conditionally closed. Two-modulus hierarchy dynamically stable. One residual gap: $B/A$ ratio not from first principles. Well-defined mathematical problem ($\sim$1--2 months).

\item \textbf{Spectral torsion}: Chamseddine et al.\ (arXiv:2511.08159) opens Yukawa correction route. 1.40\% residual mass error may be explainable.

\item \textbf{Birefringence strategy}: Conditional on $n\pi$ phase. CS evolution if $n = 0$, neutrality if $n > 0$. No action until Simons Observatory ($\sim$2027).
\end{enumerate}

\textbf{What remains genuinely open in theory}: (a) $k_{\max} = 60$ origin (anomaly cancellation approach C best but $c_1$ not computed; dropped after 7 iterations of deferral), (b) fermion mass exponents $n_f$ (spectral torsion is best lead, Grade C+), (c) $B/A$ ratio in alpha\textsuperscript{57}.
\end{tcolorbox}


%=============================================================================
\part{Agent 1: Fine Structure Constant ($\alpha$)}
\label{part:agent1}
%=============================================================================

\section{Mandate}

Maintain $\alpha^{-1} = 137.036$ derivation. Track competing derivations. Link to birefringence. \textbf{i49 PRIORITY}: Final synthesis of $\alpha$ status for programme handoff.

\section{$\alpha$ Derivation: FINAL STATUS}

\begin{tcolorbox}[colback=green!5!white, colframe=green!60!black,
  title={$\alpha$ Derivation: Definitive Status for Programme Handoff}]

\textbf{The result}:
\[
\alpha^{-1} = \frac{\pi^{3/2}}{24} \cdot \Tr(Y^2) \cdot \frac{k(k+3)}{k+4} \cdot \left[1 + \frac{7}{80(d^2 - 1)}\right] = 137.035\,999\,854
\]
with $k = k_{\max} = 60$, $d = 7$ (internal dimension), $\Tr(Y^2) = 10$ (Standard Model hypercharge sum).

\textbf{Provenance}: Derived from axioms A1--A4 + SM content. A theorem, not a fifth axiom (confirmed i40).

\textbf{Precision}: $+0.005$ ppm from CODATA 2018 value $137.035\,999\,084(21)$.

\textbf{The 0.77 ppb offset}: Physical or calculational? Three possibilities:
\begin{enumerate}
\item Two-loop correction (expected $\sim$1 ppb from higher Seeley--DeWitt coefficients).
\item Threshold correction from $\Tr(Y^2)$ running (negligible: SM is asymptotically free in $Y$).
\item Genuine prediction of a measurement that hasn't been made (the formula predicts the UV value).
\end{enumerate}
No resolution achieved in 10 iterations. Declared \textbf{OPEN but non-critical}: the formula works to sub-ppm regardless.

\textbf{Competing derivations}: Three tracked (i47--i48): $M_3(\mathbb{C})$ (philosophical, not physics), Kosmoplex (4 orders worse), tensor field (no numerical result). None credible. DFD's derivation remains unique.
\end{tcolorbox}

\section{Birefringence Accommodation: No Change from i48}

Strategy decided at i48: CS evolution (Option 3) if $n = 0$, neutrality (Option 1) if $n > 0$. The $n\pi$ phase ambiguity analysis (PRL 2026, Naokawa \& Namikawa) confirms the rotation angle may be significantly larger than $0.3^\circ$. No new data since i48. Simons Observatory ($\sim$2027) remains decisive.

\textbf{New context}: The galaxy polarization probe (PRL 134, 161001, 2025) provides a CMB-independent channel. If confirmed, this strengthens the case for an astrophysical (non-DFD) origin, favouring Option 1 regardless of $n$.

\section{Th-229: Three-Paper Foundation --- Timeline Update}

\begin{itemize}
\item Nature (Jan 2026): frequency reproducibility 220 Hz.
\item Nature Communications (2025): $K_\alpha = 5900(2300)$.
\item PRL 134, 113801 (2025): temperature sensitivity, optimal 195 K.
\end{itemize}

Path to $\dot{\alpha}/\alpha \sim 10^{-19}$/yr test by $\sim$2030. DFD predicts exactly zero. \GREEN.

\section{Agent 1 Assessment: Grade B+}

Unchanged. No new developments requiring grade change. Birefringence strategy stable. $\alpha$ derivation stable. Th-229 path advancing on schedule.

\section{New Literature (Agent 1)}

\begin{enumerate}
\item Cosmic birefringence $n\pi$ phase analysis --- continued from i47--i48.
\item Galaxy polarization probe --- continued from i47.
\item Birefringence resolves $\Sigma m_\nu$ (Namikawa mechanism) --- continued from i47.
\item PRL 134, 113801 --- Th-229 temperature sensitivity --- continued from i48.
\end{enumerate}

No genuinely new Agent 1 literature since i48. This is expected at the penultimate iteration.


%=============================================================================
\part{Agent 2: Topology and $k_{\max}$}
\label{part:agent2}
%=============================================================================

\section{Mandate}

Origin of $k_{\max} = 60$. NCG topology. \textbf{i49 PRIORITY}: Final status of $k_{\max}$ problem for handoff.

\section{$k_{\max} = 60$: FINAL STATUS}

\begin{tcolorbox}[colback=yellow!5!white, colframe=yellow!60!black,
  title={$k_{\max} = 60$: Final Status --- OPEN Problem}]

\textbf{What is known}:
\begin{itemize}
\item $k_{\max} = 60$ is a \textbf{theorem}: it is the unique value that makes the Chern--Simons level sum converge to $\alpha^{-1} = 137.036$.
\item It is \textbf{not derived from geometry alone}: it uses the SM hypercharge spectrum as input.
\item Pure gauge anomaly $\Tr(Y^3) = 0$ for all $k$ --- does NOT select $k_{\max}$ (closed i42).
\item Mixed gravitational-gauge anomaly $c_1^{(Y)}$ was the best remaining path (Approach C from i41). The representation-theoretic computation was deferred for 7 consecutive iterations (i40--i46) and formally DROPPED at i47.
\end{itemize}

\textbf{Assessment}: $k_{\max} = 60$ is an \textbf{empirically determined constant} within the DFD framework. It plays a role analogous to the number of SM generations ($N_g = 3$): unexplained but precise. The programme has failed to derive it from first principles.

\textbf{For programme handoff}: This is a defined mathematical problem. Given the $\CP^2 \times S^3$ compactification with $\text{Spin}^c$ structure, compute the mixed gravitational-gauge anomaly coefficient $c_1^{(Y)}(k)$ and determine whether it vanishes for $k = 60$ only. Estimated effort: 2--4 months of focused mathematical work by a specialist in NCG or topological field theory.
\end{tcolorbox}

\section{NCG Topology: Stable}

The $\CP^2 \times S^3$ internal space with $\text{Spin}^c$ structure remains the unique choice forced by the DFD axioms. The spectral torsion result (arXiv:2511.08159) confirms nonvanishing torsion in the internal NCG --- consistent with DFD's $S^3$ Chern--Simons sector.

\section{Agent 2 Assessment: Grade B+}

Unchanged. No new topology results. ET site selection advancing (bid book end of 2026, decision 2027) --- Agent 2 notes the EMR subsurface feasibility confirmed and Sardinia ET-SUnLab tender published.

\section{New Literature (Agent 2)}

\begin{enumerate}
\item ET-EMR: ``All signals are green'' interview (March 2026) --- subsurface allows tunnel construction, no showstoppers. \NEW.
\item ET-SUnLab: EU call for tenders published (INFN, early 2026) --- Sardinia underground laboratory construction. \NEW.
\item ET spokespersons appointed: Michele Maggiore and Ang\'elique Lartaux (March 2026). \NEW.
\item Spectral torsion of internal NCG Standard Model (arXiv:2511.08159) --- continued from i43.
\end{enumerate}


%=============================================================================
\part{Agent 3: Fermion Masses}
\label{part:agent3}
%=============================================================================

\section{Mandate}

Mass formula status. Spectral torsion corrections. $\Sigma m_\nu$ monitoring. \textbf{i49 PRIORITY}: Definitive mass formula status for handoff.

\section{Mass Formula: FINAL STATUS}

\begin{tcolorbox}[colback=green!5!white, colframe=green!60!black,
  title={Mass Formula Hierarchy: Definitive Status}]

\textbf{Established results} (Grade A or A-):
\begin{itemize}
\item VEV: $v = M_P \alpha^8 \sqrt{2\pi} = 246.09$ GeV. \textbf{Grade A.}
\item Quartic sum rule: $\sum_f m_f^4 / v^4$ matches SM content. \textbf{Grade A.}
\item Absolute mass scale from spectral action: top mass order-of-magnitude correct. \textbf{Grade A-.}
\item $\epsilon_H = 0.05$ is Higgs localization width (NOT slow-roll parameter). \textbf{Grade A.}
\end{itemize}

\textbf{Partially established} (Grade B):
\begin{itemize}
\item b-quark: $n_b = 0$ adopted as baseline (0.3\% PDG match).
\item $\Sigma m_\nu = 0.061$ eV: zero-parameter prediction. Under pressure from cosmological bounds but safe in 4 resolution paths.
\end{itemize}

\textbf{Open} (Grade C+):
\begin{itemize}
\item Individual fermion mass exponents $n_f$: not derived from first principles. Spectral torsion (arXiv:2511.08159) is the best lead but has not been worked through.
\item 1.40\% residual mass error: possibly from spectral torsion Yukawa corrections.
\end{itemize}

\textbf{Dropped items}: Mass paper Theorem K.4 edit (dropped i48 after 6 iterations of deferral).
\end{tcolorbox}

\section{$\Sigma m_\nu$ Status: Four Resolution Paths, Tightening Margin}

\begin{tcolorbox}[colback=yellow!5!white, colframe=yellow!60!black,
  title={$\Sigma m_\nu$: DFD 0.061 eV vs Cosmological Bounds --- i49 Update}]

\textbf{DFD prediction}: $\Sigma m_\nu = 0.061$ eV (zero-parameter, normal ordering).

\textbf{Current bounds}:
\begin{itemize}
\item DESI DR2 + Planck PR3 ($\Lambda$CDM): $< 0.072$ eV (95\% CL). Margin 11 meV.
\item DESI DR2 + Planck PR4 (HiLLiPoP+LoLLiPoP): relaxed, $< 0.10$ eV. Comfortable.
\item Frequentist $\Lambda$CDM: $< 0.053$ eV. \textbf{Nominally excludes DFD.}
\item $w_0 w_a$CDM: relaxed to $< 0.163$ eV. Very comfortable.
\end{itemize}

\textbf{Four resolution paths} (from i48):
\begin{enumerate}
\item Spatial curvature: relaxes $\Lambda$CDM bound from $2.6\sigma$ to $1.2\sigma$.
\item Dynamical dark energy ($w_0 w_a$CDM): relaxes to $< 0.163$ eV.
\item Birefringence (Namikawa mechanism): reionization signal absorbs part of $\Sigma m_\nu$ constraint.
\item Graham--Green--Meyers two-mass-parameter framework (PRD 113, 2026): distinguishes $M_\nu^{\text{BAO}}$ from $M_\nu^{\text{lens}}$. DFD's $\psi$-screen naturally produces this split.
\end{enumerate}

\textbf{JUNO update}: First results published (arXiv:2511.14593). World-leading $\Delta m^2_{21}$ and $\theta_{12}$ precision. Mass ordering determination requires years of additional data. Combined T2K+NOvA+JUNO: normal ordering at $2.2\sigma$. DFD's normal-ordering prediction on track.

\textbf{Assessment}: \YELLOW. Tightening but not excluded. DFD is safe in all four resolution paths. JUNO mass ordering accumulating in DFD's favour.
\end{tcolorbox}

\section{Agent 3 Assessment: Grade B+}

Unchanged from i48. Mass formula stable. $\Sigma m_\nu$ monitoring continues.

\section{New Literature (Agent 3)}

\begin{enumerate}
\item JUNO first results (arXiv:2511.14593) --- continued from i46.
\item JUNO + T2K + NOvA mass ordering --- continued from i46.
\item Graham--Green--Meyers (PRD 113, 043514) --- continued from i48.
\item DESI DR2 $\Sigma m_\nu$ bounds --- continued from i48.
\item Scalar NSI risk to JUNO (arXiv:2602.05564) --- continued from i46.
\item arXiv:2603.17181 (March 2026) --- Impact of new physics on JUNO--long-baseline synergy. \NEW.
\end{enumerate}


%=============================================================================
\part{Agent 4: Spin$^c$ and Exponents}
\label{part:agent4}
%=============================================================================

\section{Mandate}

$\text{Spin}^c$ structure of $\CP^2 \times S^3$. Phase 0B derivation of perturbation parameters. \textbf{i49 PRIORITY}: What must Phase 0A show for theory to proceed.

\section{What Phase 0A Must Show: Theory Requirements}

\begin{tcolorbox}[colback=red!5!white, colframe=red!60!black,
  title={CRITICAL: What Phase 0A Must Demonstrate for Theory}]

Phase 0A (background-only SymBoltz.jl run) must answer:

\begin{enumerate}
\item \textbf{CMB power spectrum with DFD $G_{\text{eff}}(k,a)$}: Does the deep-MOND $G_{\text{eff}} = G_N \sqrt{a_0 / (4\pi G \rho_b \delta_b a)}$ produce a viable $C_\ell$ spectrum?

\item \textbf{Third-peak ratio $R_{31}$}: Analytic estimates range from 0.41 to 0.70 vs observed $\sim$1.0. This is a 10--61\% deficit. Only numerical computation can settle it. If $R_{31} < 0.5$: DFD requires a vector-field extension (AeST-like). If $R_{31} \geq 0.7$: DFD survives as-is.

\item \textbf{$\sigma_8$ numerical confirmation}: Analytic estimate $\sigma_8 = 0.83$ (+0.22/-0.08). Phase 0A gives the exact number.

\item \textbf{$A_L^{\text{DFD}}$}: Analytic revised to 1.00--1.05 (i42). Current data: $1.025 \pm 0.017$. Consistent but needs numerical confirmation.

\item \textbf{BAO amplitude}: DFD predicts enhanced BAO amplitude ($\times 1.3$--$2.0$). Phase 0A gives precise number.
\end{enumerate}

\textbf{Go/no-go criterion}: If $R_{31} < 0.4$, the cosmological sector of DFD is dead. If $R_{31} > 0.6$, DFD is alive. Between 0.4 and 0.6: marginal, requires Phase 0B (perturbation theory).

\textbf{Timeline}: SymBoltz.jl is published (A\&A, March 2026). Phase 0A requires 12 hours of Julia computation. This is NOT a theory problem; it is a computation-execution problem.
\end{tcolorbox}

\section{Phase 0B: Perturbation Theory Parameters}

Five parameters underived (from i45):
\begin{enumerate}
\item $c_s$ at cosmological epochs (confirmed $\sim c$; i18 established, i20 refined).
\item $G_{\text{eff}}(k,a)$ functional form (deep-MOND quasi-static; i20--i21).
\item Screening function form (Appendix Q vs derived; i23 gap).
\item Silk damping scale interaction with $G_{\text{eff}}$.
\item Initial conditions at recombination.
\end{enumerate}

Phase 0B begins only after Phase 0A succeeds. Estimated: weeks--months.

\section{Agent 4 Assessment: Grade B}

Unchanged. Phase 0B framework outlined but Phase 0A is the blocker. 8th iteration waiting.

\section{New Literature (Agent 4)}

No new Agent 4 literature since i48.


%=============================================================================
\part{Agent 5: Anomaly and $k$-Selection}
\label{part:agent5}
%=============================================================================

\section{Mandate}

Mixed anomaly $c_1^{(Y)}$. $\Sigma m_\nu$ context. \textbf{i49 PRIORITY}: Final status.

\section{Mixed Anomaly: DROPPED (Final)}

Mixed gravitational-gauge anomaly $c_1^{(Y)}$ computation was deferred for 7 iterations (i40--i46) and formally dropped at i47. The $k_{\max} = 60$ origin remains an open mathematical problem for the programme handoff document. See Agent 2 for the defined problem statement.

\section{$\Sigma m_\nu$: Sterile Neutrino and NSI Context}

\begin{itemize}
\item arXiv:2603.17181 (March 2026) examines JUNO--long-baseline synergy under new physics scenarios including sterile neutrinos and NSI. DFD predicts NO sterile neutrinos (zero additional neutrino mass states). If JUNO+T2K+NOvA find evidence for a sterile neutrino, DFD's $\Sigma m_\nu = 0.061$ eV prediction survives (it refers only to active neutrinos) but the cosmological interpretation changes.
\item Scalar NSI (arXiv:2602.05564) remains a risk to JUNO's mass ordering sensitivity but does not affect DFD's prediction.
\end{itemize}

\section{Agent 5 Assessment: Grade C+}

Unchanged. Limited scope after mixed anomaly dropped.


%=============================================================================
\part{Theory Sector: Programme Handoff Summary}
\label{part:handoff}
%=============================================================================

\section{Defined Open Problems for External Researchers}

\begin{center}
\begin{longtable}{p{3cm}p{6cm}p{2cm}p{2.5cm}}
\toprule
\textbf{Problem} & \textbf{Description} & \textbf{Effort} & \textbf{Skills Needed} \\
\midrule
\endfirsthead
$k_{\max} = 60$ & Compute mixed anomaly $c_1^{(Y)}(k)$ on $\CP^2 \times S^3$. Determine if it vanishes uniquely at $k = 60$. & 2--4 months & NCG, TFT \\
Fermion exponents $n_f$ & Derive individual fermion mass power-law exponents from spectral torsion of internal NCG & 3--6 months & NCG, spectral geometry \\
$B/A$ ratio & Complete alpha\textsuperscript{57} Step 9 by deriving the modulus ratio from first principles & 1--2 months & Algebraic geometry \\
0.77 ppb offset & Determine if two-loop Seeley--DeWitt corrections account for the offset & 2--3 months & Heat kernel, QFT \\
\bottomrule
\end{longtable}
\end{center}

\section{Theory Sector Honest Assessment}

\begin{tcolorbox}[colback=blue!5!white, colframe=blue!60!black,
  title={Theory Sector: Honest Final Assessment}]

\textbf{Achievements}:
\begin{itemize}
\item Sub-ppm $\alpha$ derivation from geometry alone (unprecedented in theoretical physics).
\item VEV formula $v = M_P \alpha^8 \sqrt{2\pi}$ confirmed.
\item 3-axiom minimal structure with 21:1 prediction-to-parameter ratio.
\item Two-component decomposition ($\psi = \psi_{\text{grav}} + \delta\psi_{\text{EM}}$) resolved all 5 original tensions (T1--T5).
\item Distance-bias framework (i14) correctly identifies DFD's observational signature.
\end{itemize}

\textbf{Failures}:
\begin{itemize}
\item $k_{\max} = 60$ origin never derived (7 iterations of deferral, then dropped).
\item Individual fermion mass exponents not derived.
\item Mixed anomaly never computed.
\item Phase 0A (numerical Boltzmann computation) never executed in 8 iterations.
\item Mass paper Theorem K.4 edit never completed (6 iterations, then dropped).
\end{itemize}

\textbf{Honest grade for theory sector}: \textbf{B+}. The achievements are genuine and significant. The failures are real and documented. The framework is internally consistent and makes testable predictions. The open problems are well-defined and tractable by specialists.
\end{tcolorbox}


\end{document}
